\documentclass[11pt]{article}

\def\chapitre{15}
\def\pagetitle{Endomorphismes des espaces euclidiens.}

\input{setup.tex}

\begin{document}

\input{title.tex}

\renewcommand*{\D}{\cursive{D}}
\newcommand*{\SO}{\nt{SO}}

\section{Généralités.}

Ici, $(E,\Lg\.,\.\Rg)$ désignera toujours un espace euclidien.

\begin{defi}{}{}
    Soit $u\in\L(E)$, $u$ est une \bf{isométrie} (ou endomorphisme orthogonal) ssi $\|u(x)\|=\|x\|$.
\end{defi}

\vspace*{-0.8cm}

\begin{prop}{}{}
    Une isométrie est bijective.
    \tcblower
    Soit $u$ une isométrie, et $x\in\Ker(u)$, alors $\|x\|=\|u(x)\|=0_\R$ donc $x=0_E$. Donc $u$ est injective.\\
    Or $\dim E \leq \infty$, et $u\in\L(E)$, donc $u$ injective ssi $u$ bijective.\\
    \bf{Remarque.} La seule projection bijective est l'identité.
\end{prop}

\vspace*{-0.8cm}

\begin{nota}{}{}
    Notons $O(E)$ l'ensemble des isométries de $E$ et $SO(E)$ l'ensemble des éléments de $O(E)$ de déterminant $1$.
\end{nota}

\vspace*{-0.8cm}

\begin{prop}{}{}
    \begin{enumerate}
        \item $(O(E),\circ)$ est un groupe.
        \item $(SO(E),\circ)$ est un sous-groupe de $(O(E),\circ)$.
    \end{enumerate}
    \tcblower
    \boxed{1.} On a $\Id_E\in O(E)$ donc $O(E)$ non-vide; stable par produit :\\
    $\forall u,v \in O(E)$, $\forall x \in E, ~ \|(u\circ v)(x)\|=\|u(v(x))\|=\|v(x)\|=\|x\|$ donc $u\circ v \in O(E)$.\\
    Stable par inverse : pour $u\in O(E)$, $\forall x \in E, ~ \|u^{-1}(x)\|=\|u(u^{-1}(x))\|=\|x\|$ donc $u^{-1}\in O(E)$.\\
    \boxed{2.} Facile.
\end{prop}

\begin{thm}{Caractérisation des isométries. (1)}{}
    Soit $u\in\L(E)$. Les deux propositions suivantes sont équivalentes :
    \begin{enumerate}
        \item $\forall x \in E, ~ \|u(x)\|=\|x\|$.
        \item $\forall (x,y) \in E^2, ~ \Lg u(x), u(y)\Rg = \Lg x,y\Rg$.
    \end{enumerate}
    \tcblower
    \boxed{2 \ra 1} Avec $y=x$ : $\|u(x)\|=\|x\|$.\\
    \boxed{1\ra 2} Soit $(x,y)\in E^2$, alors $\|u(x+y)\|=\|x+y\|$ puis $\|u(x+y)\|=\|u(x)+u(y)\|$.\\
    Donc $\|x+y\|^2=\|u(x)+u(y)\||^2$ donc $\|x\|^2+2\Lg x,y\Rg + \|y\|^2=\|u(x)\|^2+2\Lg u(x),u(y)\Rg+\|y\|^2$.\\
    De plus, $\|x\|^2=\|u(x)\|^2$ et $\|y\|^2=\|u(y)\|^2$ donc $2\Lg x,y\Rg = 2\Lg u(x), u(y)\Rg$.\\
    Finalement, $\Lg x,y \Rg = \Lg u(x), u(y)\Rg$.\\
    \bf{Remarque.} Pour $x,y\in E$ d'écart angulaire $\theta$ : $\cos\theta=\frac{\Lg x,y \Rg}{\|x\|\|y\|}$ alors $\frac{\Lg u(x),u(y)\Rg}{\|u(x)\|\|u(y)\|}=\frac{\Lg x,y\Rg}{\|x\|\|y\|}$.
\end{thm}

\vspace*{-0.8cm}

\begin{thm}{Caractérisation des isométries. (2)}{}
    Soit $u\in \L(E)$. Alors $u$ est une isométrie ssi $u^* = u^{-1}$ ssi $u\circ u^* = u^* \circ u = \Id_E$.
    \tcblower
    D'après la proposition précédente, $u\in O(E)\iff \forall (x,y) \in E^2, ~ \Lg u(x),u(y)\Rg=\Lg x,y\Rg$.\\
    Or par définition de l'adjoint, $\forall (x,y)\in E^2, ~ \Lg x,u(y)\Rg = \Lg u^*(x),y\Rg$.\\
    Ainsi, $u\in O(E) \iff \forall (x,y)\in E^2, ~ \Lg u^*(u(x)), y \Rg = \Lg x,y\Rg \iff \forall x \in E, ~ u^*(u(x))=x \iff u^*\circ u = \Id_E$.
\end{thm}

\vspace*{-0.8cm}

\bf{Remarque.} Les isométries font partie des endomorphismes normaux (HP) : $u\in\L(E)$ normal $\iff$ $u\circ u^* = u^* \circ u$.

\begin{thm}{Caractérisation des isométries. (3)}{}
    Soit $u\in\L(E)$. Les proposition suivantes sont équivalentes :
    \begin{enumerate}
        \item $u$ est une isométrie de $E$.
        \item L'image par $u$ de n'importe quelle b.o.n. de $E$ est une b.o.n. de $E$.
        \item L'image par $u$ d'une b.o.n. de $E$ est une b.o.n. de $E$.
    \end{enumerate}
    \tcblower
    \boxed{1\ra 2} Pour $u\in O(E)$, et $\B=(e_1,...,e_n)$ b.o.n. de $E$, $\forall i \in \lb1,n\rb, ~ \|u(e_i)\|=\|e_i\|=1$.\\
    Puis $\forall i,j \in \lb1,n\rb, ~ i\neq j \ra \Lg u(e_i), u(e_j)\Rg = \Lg e_i, e_j \Rg = 0$.\\
    Donc $\B'=(u(e_1),...,u(e_n))$ est une famille orthonormale de $E$ et $|\B'|=\dim E < \infty$ donc $\B'$ est b.o.n de $E$.\\
    \boxed{2\ra3} Trivial.\\ 
    \boxed{3\ra 1} Pour $\B=(e_1,...,e_n)$ b.o.n. de $E$ tq $u(\B)=(e'_1,...,e'_n)$ soit encore b.o.n. de $E$.\\
    Soit $x\in E$ : $x=\sum_{i=1}^nx_ie_i$, puis $\|x\|^2=\sum_{i=1}^nx_i^2$ car b.o.n. et $u(x)=\sum_{i=1}^nx_iu(e_i)=\sum_{i=1}^nx_ie'_i$.\\
    Puis $\|u(x)\|^2=\sum_{i=1}^nx_i^2$ car $\B'$ est b.o.n. puis $\forall x \in E, ~ \|x\|=\|u(x)\|$ donc $u\in O(E)$.
\end{thm}

\vspace*{-0.8cm}

\begin{thm}{Caractérisation des isométries. (4)}{}
    Soit $u\in\L(E)$. Les propositions suivantes sont équivalentes :
    \begin{enumerate}
        \item $u\in O(E)$.
        \item Soit $\B$ une b.o.n. de $E$, alors $\Mat_\B(u)\in O_n(\R)$.
        \item Il existe $\B$ une b.o.n. de $E$ telle que $\Mat_\B(u)\in O_n(\R)$.
    \end{enumerate}
    \tcblower
    \boxed{1\ra2} $u\circ u^*=\Id_E \iff \Mat_\B(u)\Mat_\B(u)^\top=I_n$ i.e. $\Mat_\B(u)\in O_n(\R)$.\\
    \boxed{2\ra3} Trivial.\\
    \boxed{3\ra 1} $\Mat_\B(u)\Mat_\B(u)^\top = I_n \ra u\circ u^* = \Id_E$.
\end{thm}

\begin{rappel}{}{}
    Les propriétés démontrées sur les isométries sont les mêmes que celles sur les matrices orthogonales.
    \begin{equation*}
        \begin{array}{|c|c|}
            \hline
            u \in O(E) & \Mat_\B(u) \in O_n(\R)\\
            \hline
            u $ est inversible et $ u^{-1}=u^* & A $ est inversible et $ A^{-1} = A^\top\\
            \hline
            \forall x \in E, ~ \|u(x)\|=\|x\| & \forall X \in \M_{n,1}(\R), ~ \|AX\|=\|X\|\\
            \hline
            \forall x, y \in E, ~ \Lg u(x), u(y) \Rg = \Lg x,y\Rg & \forall X,Y \in \M_{n,1}(\R), ~ \Lg AX, AY \Rg = \Lg X,Y\Rg\\
            \hline
            u $ envoie une b.o.n. sur une b.o.n.$ & $Les colonnes de $ A $ sont orthonormées par le produit scalaire usuel de $ \R^n.\\
            \hline
            \Sp(u) \subset \{-1,1\} & \Sp A \subset \{-1,1\}\\
            \hline 
        \end{array}
    \end{equation*}
\end{rappel}

\begin{prop}{}{}
    Soit $u\in O(E)$, alors $\det(u)\in\{-1,1\}$.
    \tcblower
    Soit $\B$ une b.o.n. de $E$ et $u\in O(E)$, alors $\Mat_\B(u)\in O_n(\R)$ donc $\det(u)=\det(A)\in\{-1,1\}$.
\end{prop}

\begin{prop}{}{}
    Soit $s$ une symétrie de $E$, alors $s$ est une isométrie ssi $s$ est une symétrie orthogonale.
    \tcblower
    \boxed{\la} Soit $F=\Ker(s-\Id_E)$ et $x\in E$ : $\exists y\in F, ~ \exists z\in F^\bot, ~ x=y+z$ puis $s(x)=y-z$.\\
    Ainsi, $\|x\|^2=\|y\|^2+\|z\|^2=\|s(x)\|^2$ avec le théorème de Pythagore.\\
    \boxed{\ra} Soient $F=\Ker(s-\Id_E)$ et $G=\Ker(s+\Id_E)$. Il s'agit de démontrer que $F\bot G$.\\
    Soit $x\in F$ et $y\in G$. Alors $s(x+y)=x-y$ et $\|x-y\|=\|s(x+y)\|=\|x+y\|$.\\
    Donc $\|x-y\|^2=\|x+y\|^2$ donc $\|x\|^2-2\Lg x,y\Rg + \|y\|^2 = \|x\|^2 + 2\Lg x,y \Rg + \|y\|^2$ donc $\Lg x,y \Rg = 0$.
\end{prop}

\begin{ex}{}{}
    Soit $u\in\L(E)$ tel que $\forall (x,y) \in E^2, ~ \Lg x,y \Rg = 0 \ra \Lg u(x), u(y) \Rg = 0$.\\
    Montrer qu'il existe $\l\in\R^*, ~ \forall x \in E, ~ \|u(x)\|=\l\|x\|$.
    \tcblower
    \bf{Analyse.} Supposons que $\l$ existe, alors pour $x_0\in S(0,1)$ : $\|u(x_0)\|=\l\|x_0\|=\l$.\\
    Soient $x,y \in S(0,1)$. $\Lg x+y, x-y \Rg = \|x\|^2+\Lg y,x \Rg - \Lg x,y\Rg - \|y\|^2 = 0$.\\
    Donc $\Lg u(x+y), u(x-y) \Rg = 0$ donc $\Lg u(x) + u(y), u(x) - u(y) \Rg = 0$ donc $\|u(x)\|^2=\|u(y)\|^2$.\\
    Soit $x\in E$, si $x=0$, alors $u(x)=\l x$ est vérifié.\\
    Sinon, $\l = \|u(\frac{x}{\|x\|})\|=\frac{1}{\|x\|}\|u(x)\|$, on a bien $\|u(x)\|=\l\|x\|$.
\end{ex}

\section{Isométries du plan.}

\begin{prop}{}{11}
    \begin{equation*}
        \SO_2(\R) = \left\{ \begin{pmatrix}\cos\theta&-\sin\theta\\\sin\theta&\cos\theta \end{pmatrix}, ~ \theta\in\R \right\}; \quad R(\theta) = \begin{pmatrix}\cos\theta&-\sin\theta\\\sin\theta&\cos\theta\end{pmatrix}.
    \end{equation*}
    Alors $\forall \theta,\theta'\in\R, ~ R(\theta+\theta')=R(\theta)\times R(\theta')$.
    \tcblower
    \boxed{\subset} Soit $A\in\SO_2(\R)$.
    \begin{equation*}
        A=\begin{pmatrix}a&b\\c&d\end{pmatrix}; \quad \begin{cases}a^2+c^2 = 1\\b^2+d^2=1\\ab+cd=0\end{cases}
    \end{equation*}
    On a $\det(A) = 1 \iff ad-bc = 1$ et $a^2+c^2=1 \ra |a|\leq 1 \ra \exists\theta\in\R, ~ a=\cos\theta$ puis $c=\sin\theta$.\\
    De même, $d^2+b^2=1$ donc $\exists \theta'\in\R, ~ d=\cos\theta'$ et $b=-\sin\theta'$. On peut supposer $\theta,\theta'\in[0,2\pi[$.\\
    Il reste $ab+cd = 0$ et $ad-bc=1$ doncc $\sin(\theta-\theta')=0$ et $\cos(\theta-\theta')=1$ donc $\theta=\theta'$.\\
    \boxed{\supset} $R(\theta)\in O_2(\R)$ et $\det R(\theta)=\cos^2\theta+\sin^2\theta=1$ donc $R(\theta)\in \SO_2(\R)$.\\
    Soient $\theta,\theta'\in\R$.
    \begin{equation*}
        R(\theta)R(\theta')=\begin{pmatrix}\cos\theta\cos\theta'-\sin\theta\sin\theta' & -\cos\theta\sin\theta'-\sin\theta\cos\theta'\\\sin\theta\cos\theta'+\cos\theta\sin\theta'&\cos\theta\cos\theta'-\sin\theta\sin\theta'\end{pmatrix} = \begin{pmatrix}\cos(\theta+\theta')&-\sin(\theta+\theta')\\\sin(\theta+\theta')&\cos(\theta+\theta')\end{pmatrix}=R(\theta+\theta').
    \end{equation*}
\end{prop}

\begin{exercice}{}{}
    Soit $R(\theta)\in\SO_2(\R)$ avec $\theta\in[0,2\pi[$.
    \begin{enumerate}
        \item Montrer que $R(\theta)$ est diagonalisable ssi $\theta\in\{0,\pi\}$.
        \item Montrer que $R(\theta)\in\M_2(\C)$ est toujours diagonalisable.
    \end{enumerate}
    \tcblower
    \boxed{1.} $\X_{R(\theta)}=X^2-\Tr(R(\theta))+\det(R(\theta))=X^2-2\cos\theta X+1$, $\Delta=4(\cos^2\theta-1)\leq 0$.\\
    Si $\theta\in]0,2\pi[$ et $\theta\neq\pi$, alors $\Delta<0$ et $\X_{R(\theta)}$ n'est pas scindé donc pas diagonalisable.\\
    Si $\theta=0$, alors $R(\theta)=I_n$ diagonalisable. Si $\theta=\pi$, alors $R(\theta)=-I_2$ diagonalisable.\\
    \boxed{2.} On a toujours $\X_{R(\theta)}=X^2-2\cos\theta X+1 = (X-e^{i\theta})(X-e^{-i\theta})$.\\
    Si $\theta\in\{0,\pi\}$, alors $R(\theta)$ est diagonalisable dans $\M_2(\R)$ donc dans $\M_2(\C)$.\\
    Sinon, $e^{i\theta}\neq e^{-i\theta}$ puis $\X_{R(\theta)}$ est SRS dans $\C[X]$, donc $R(\theta)$ est diagonalisable.
\end{exercice}

\begin{prop}{Rotation dans un espace orienté.}{}
    Soit $(E, \Lg\.,\.\Rg)$ euclidien de dimension 2 et orienté.\\
    Soit $u\in\SO(E)$ et $\B$ une base directe de $E$. Alors $\exists!\theta\in[0,2\pi[, ~ \Mat_\B(u)=R(\theta)$.
    \tcblower
    Puisque $\B$ est une b.o.n. de $E$, $\Mat_\B(u)\in O_2(\R)$, et $\det(u)=1$ donc $\det(\Mat_\B(u))=1$.\\
    Donc $\exists!\theta\in[0,2\pi], ~ \Mat_\B(u)=R(\theta)$.\\
    Soient $\B,\B'$ deux bases directes de $E$. Notons $P\in \SO_2(\R)$ la matrice de passage de $\B$ à $\B'$.\\
    Alors $\Mat_\B(u)=P\Mat_{\B'}(u)P^{-1}$, donc $\exists!\theta''\in[0,2\pi[, ~ P=R(\theta'')$.\\
    Donc $\Mat_\B(u)=P\Mat_{\B'}(u)P^{-1} \iff R(\theta)=R(\theta'')R(\theta')(R(\theta)'')^{-1}=R(\theta')$.
\end{prop}

\begin{prop}{}{}
    \begin{equation*}
        O_2(\R) \setminus \SO_2(\R) = \left\{ \begin{pmatrix}\cos\theta & \sin\theta \\ \sin\theta & -\cos\theta \end{pmatrix}, ~ \theta\in[0,2\pi[ \right\}.
    \end{equation*}
    \tcblower
    \boxed{\subset} Pareil qu'en proposition \ref{prop:11}.\\
    \boxed{\supset} Soit $S(\theta)=\begin{pmatrix}\cos\theta&\sin\theta \\ \sin\theta & -\cos\theta\end{pmatrix}$.\\
    On a $\det(S(\theta))=-(\cos^2\theta+\sin^2(\theta))=-1$ donc $S(\theta)\in O_2(\R)\setminus\SO_2(\R)$, et $S(\theta)^2=I_2$.
\end{prop}

\begin{prop}{}{}
    $S(\theta)=\begin{pmatrix}\cos\theta&\sin\theta\\\sin\theta&-\cos\theta\end{pmatrix}\in O_2(\R)\setminus\SO_2(\R)$ est diagonalisable pour tout $\theta\in\R$ et $S(\theta)=R(\frac{\theta}{2})\begin{pmatrix}1&0\\0&-1\end{pmatrix}R(\frac{\theta}{2})^{-1}$.
    \tcblower
    $\X_{S(\theta)}=X^2-\Tr(S(\theta))X+\det(S(\theta))=(X-1)(X+1)$, donc $\X_{S(\theta)}$ est SRS donc $S(\theta)$ est diagonalisable.\\
    De plus, par Cayley-Hamilton, $S(\theta)^2=I_n$ donc $S(\theta)$ est une symétrie.
    \begin{align*}
        R\left( \frac{\theta}{2} \right)\begin{pmatrix}1&0\\0&-1\end{pmatrix}R\left( \frac{\theta}{2} \right)^{-1} &= \begin{pmatrix}\cos\frac{\theta}{2} & \sin\frac{\theta}{2} \\ \sin\frac{\theta}{2} & -\cos\frac{\theta}{2} \end{pmatrix} \begin{pmatrix}\cos\frac{\theta}{2} & \sin\frac{\theta}{2} \\ -\sin\frac{\theta}{2} & \cos\frac{\theta}{2}\end{pmatrix} =\begin{pmatrix}\cos^2\frac{\theta}{2}-\sin^2\frac{\theta}{2} & 2\sin\frac{\theta}{2}\cos\frac{\theta}{2} \\ 2\sin\frac{\theta}{2}\cos\frac{\theta}{2} & -\cos^2\frac{\theta}{2}+\sin^2\frac{\theta}{2}\end{pmatrix}\\
        &=\begin{pmatrix}\cos\theta & \sin\theta \\ \sin\theta & -\cos\theta \end{pmatrix} = S(\theta).
    \end{align*}
\end{prop}

\begin{prop}{}{}
    Soit $E$ euclidien de dimension 2 et $u\in O(E)\setminus\SO(E)$, $u$ est une symétrie orthogonale par rapport à une droite.
    \tcblower
    Soit $\B$ une b.o.n. de $E$, $u$ est isométrie et $\B$ b.o.n. donc $\Mat_\B(u)\in O_2(\R)$, et $\det u=-1$.\\
    Donc $\exists \theta \in [0,2\pi[, ~ \Mat_\B(u)=S(\theta)$. On a vu $A^2=I_2$ donc $u^2=\Id_E$, c'est une symétrie.\\
    Or une symétrie est une isométrie ssi c'est une symétrie orthogonale, donc $u$ est une symétrie orthogonale.\\
    Puis $\det(u)=-1$, donc $u$ n'est pas triviale, donc $u$ est symétrie orthogonale par rapport à $D=\Ker(u-\Id_E)$.
\end{prop}

\begin{ex}{}{}
    Soit $E$ euclidien de dimension 2, $u\in\SO(E)$ et $a\in E$ unitaire. On note $\theta$ la mesure d'angle orienté de $u$.\\
    $\hra$ $\cos\theta=\Lg a, u(a) \Rg$ et $\sin\theta=[a, u(a)]$.
    \tcblower
    Soit $\B=(e_1,e_2)$ une b.o.n. directe quelconque de $E$. Alors $a=a_1e_1+a_2e_2$; $\Mat_\B(u)=R(\theta)$.\\
    Alors $\Lg a, u(a) \Rg  = a_1(\cos\theta a_1 - \sin\theta a_2) + a_2(\sin\theta a_1 + \cos\theta a_2) = \cos\theta\|a\|^2=\cos\theta$.\\
    Et $[a,u(a)]=\det_\B(a,u(a))=\sin(\theta)(a_1^2+a_2^2)=\sin\theta$.
\end{ex}

\section{Réduction des isométries.}

\begin{prop}{}{}
    Soit $u\in O(E)$ et $F$ un sous-espace de $E$ stable par $u$. Alors $u(F) = F$ et $u(F^\bot)\subset F^\bot$.\\
    De plus, $u_F\in O(F)$ et $u_{F^\bot}\in O(F^\bot)$.
    \tcblower
    Notons $n$ la dimension de $E$ et $p$ la dimension de $F$.\\
    Par hypothèse, $u(F)\subset F$, et $u$ est bijectif donc $\dim u(F) = \dim F$ donc $u(F)=F$.\\
    Soit $y\in u(F^\bot)$ : $\exists x \in F^\bot, ~ y=u(x)$ et $z\in F$ : $\exists t \in F, ~ z=u(t)$.\\
    Puis $\Lg y, z \Rg = \Lg u(x), u(t) \Rg = \Lg x,t\Rg = 0$ donc $u(F^\bot)\subset F^\bot$.\\
    On a $\forall x \in E, ~ \|u(x)\|=\|x\| \ra \forall x \in F, ~ \|u_F(x)\|=\|x\|$, les endomorphismes induits sont toujours isométries.
\end{prop}

\vspace*{-0.5cm}

\begin{rappel}{Chapitre 9, Exemple 22.}{}
    Soit $E$ un $\R$-evdf et $u\in\L(E)$. Alors il existe au moins une droite ou un plan stable par $u$.
\end{rappel}

\vspace*{-0.5cm}

\begin{thm}{Réduction des isométries. $\star$}{}
    \begin{enumerate}
        \item Soit $E$ de dimension $n$ et $u\in O(E)$, alors il existe $\B$ telle que $\Mat_\B(u)=\Diag(R(\theta_1),...,R(\theta_s),I_p,-I_q)$.\\
        Avec $2s+p+q = n$.
        \item Soit $A\in O_n(\R)$. Alors $A$ est orthogonalement semblable à une matrice de la forme précédente.
    \end{enumerate}
    \tcblower
    \boxed{1.} Par récurrence forte sur $n$.\\
    Si $n=1$, alors les seules applications linéaires en dimension 1 sont les homothéties.\\
    Les seules homothéties qui sont des isométries sont $\pm\Id_E$, le résultat est trivial.\\
    Si $n=2$, alors $u$ est une isométrie du plan :\\
    $\hra$ Si $u$ est une rotation, il existe $\B$ b.o.n. de $E$ et $\theta\in\R$ tq $\Mat_\B(u)=R(\theta)$.\\
    $\hra$ Si $u$ est une symétrie orthogonale selon une droite, il existe $\B$ b.o.n. de $E$ telle que $\Mat_\B(u)=\Diag(1,-1)$.\\
    Soit $n\geq3$ tel que $\forall k < n$ ... Soit $u\in O(E)$ et $H$ stable par $u$, alors $u(H^\bot)\subset H^\bot$.\\
    On a $\dim(H)<\dim(E)=n+1$ et $\dim(H^\bot)\leq \dim(E)$, on applique l'hypothèse de récurrence.\\
    Il existe $\B_1$ base de $H$ et $\B_2$ base de $H^\bot$ tq $\Mat_{\B_1}(u_{|H})$ et $\Mat_{\B_2}(u_{|H^\bot})$ de la bonne forme.\\
    On pose $\B=\B_1\sqcup \B_2$ qui est b.o.n. de $E$, et $\Mat_\B(u)=\Diag(\Mat_{\B_1}(u_{|H}), \Mat_{\B_2}(u_{|H^\bot}))$.\\
    Quitte à permuter les vecteurs de $\B$, $\Mat_\B(u)$ a la forme voulue.\\
    \boxed{2.} Soit $A\in O_n(\R)$, et $u$ son alca. D'après le 1, il existe $\B$ b.o.n. de $E$ telle que $\Mat_\B(u)=...$\\
    Notons $P$ la matrice de passage de $\B_C$ à $\B$, alors $P\in O_n(\R)$ puis $A=P\Diag(...)P^{-1}$.\\
    \bf{Remarque.} On peut imposer l'unicité avec $p,q\in\{0,1\}$.
\end{thm}

\vspace*{-0.5cm}

\begin{corr}{Réduction des matrices de rotation.}{}
    Soit $A\in\SO_n(\R)$.\\
    $\hra$ Si $n$ est pair, $A$ est orthogonalement semblable à $\Diag(R(\theta_1),...,R(\theta_s))$ avec $2s=n$.\\
    $\hra$ Si $n$ est impair, $A$ est orthogonalement semblable à $\Diag(R(\theta_1),...,R(\theta_s),1)$.
    \tcblower
    Par le théorème précédent, $A$ est orthogonalement semblable à $\Diag(R(\theta_1),...,R(\theta_s),I_p,-I_q)$.\\
    Et $\det(A)=(-1)^q$ donc $\det(A)=1\iff q$ pair donc $q=0$ et si $n$ pair, $p=n-2s=0$; sinon $p=1$.
\end{corr}

\begin{prop}{}{}
    $\SO_n(\R)$ est connexe par arcs.
    \tcblower
    Soit $A\in\SO_n(\R)$. Par théorème de réduction, il existe $P\in O_n(\R)$, ~ $A=P\Diag(R(\theta_1),...,R(\theta_s),I_p)P^{-1}$.\\
    On pose $\phi:t\mapsto P\Diag(R(t\theta_1),...,R(t\theta_s),I_p)P^{-1}\in\SO_n(\R)$, continue sur $[0,1]$ car tous les coefficients de $\phi(t)$ sont des combinaisons linéaires de $\cos(t\theta_i)$ et $\sin(t\theta_i)$, et $\phi(0)=I_n$, $\phi(1)=A$.
\end{prop}

\begin{ex}{}{}
    \begin{enumerate}
        \item Déterminer les matrices de $O_n(\R)$ diagonalisables.
        \item Montrer que toutes les matrices de $O_n(\R)$ sont diagonalisables dans $\M_n(\C)$.
    \end{enumerate}
    \tcblower
    \boxed{1.} Soit $A\in O_n(\R)$ : $A\sim \Diag(R(\theta_1),...,R(\theta_s),I_p,-I_q)$ avec $2s+p+q=n$.\\
    Puis $A$ diagonalisable ssi $\forall i \in \lb1,s\rb, ~ R(\theta_i)$ diagonalisable ssi $s=0$ ssi $A\sim\Diag(1,...,1,-1,...,-1)$ ssi $A^2=I_n$.\\
    Donc $A$ est diagonalisable ssi son endomorphisme canoniquement associé est une symétrie orthogonale.\\
    \boxed{2.} On a vu que $\forall \theta \in\R, ~ R(\theta)$ est diagonalisable dans $\M_2(\C)$.\\
    Donc $\Diag(R(\theta_1),...,R(\theta_s),I_p,-I_q)$ est diagonalisable dans $\M_n(\C)$ donc $A$ aussi.
\end{ex}

\begin{prop}{Cas de la dimension 3.}{}
    Soit $E$ euclidien de dimension 3, et $u$ une rotation de $E$ non triviale.\\
    Alors $1\in\Sp(u)$ et $\dim E_1 = 1$, où $\D=E_1$ est appelé \bf{axe de la rotation}.\\
    Alors $\cursive{H}=\D^\bot$ est un plan stable par $u$, $u_{|\cursive{H}}$ est une rotation du plan $\cursive{H}$.\\
    Il existe $\B$ une b.o.n. directe de $E$ telle que $\Mat_\B(u)=\Diag(R(\theta),1)$.
    \tcblower
    D'après le théorème de réduction des isométries, il existe $\B=(e_1,e_2,e_3)$ b.o.n. de $E$ tq $\Mat_\B(u)=\Diag(R(\theta),1)$.\\
    Alors $u(e_3)=e_3$ donc $1\in\Sp(u)$, puis $\dim E_1=1$ donc $\D=\Vect(e_3)$ est axe de la rotation.\\
    Puis $\cursive{H}=\D^\bot$ est stable par $u$ et $u_{|\cursive{H}}\in O(\cursive{H})$ isométrie du plan non triviale.\\
    Si $e_3$ est choisi comme vecteur directeur de $\D$, $\B=(e_1,e_2,e_3)$ est directe. Sinon $\B=(e_2,e_1,e_3)$ l'est. 
\end{prop}

\begin{exercice}{}{}
    Soit $E$ euclidien de dimension 3 et $u\in O(E) \setminus \SO(E)$ isométrie négative.
    \begin{enumerate}
        \item Montrer qu'il existe $\B$ une b.o.n. de $E$ telle que $\Mat_\B(u)=\Diag(R(\theta),-1)$.
        \item Montrer que $u$ est la composée d'une rotation et d'une symétrie orthogonale. On pourra calculer :
        \begin{equation*}
            \begin{pmatrix}\cos\theta&-\sin\theta&0\\\sin\theta&\cos\theta&0\\0&0&1\end{pmatrix}\begin{pmatrix}1&0&0\\0&1&0\\0&0&-1\end{pmatrix}
        \end{equation*}
        \item Vérifier que $\Tr(u)=2\cos\theta-1$.
        \item Étudier les cas particuliers  : $\theta=0$ et $\theta=\pi$.
    \end{enumerate}
\end{exercice}

\section{Endomorphismes symétries (et antisymétriques).}

\subsection{Généralités.}

\begin{defi}{}{}
    Soit $(E,\Lg\.,\.\Rg)$ euclidien et $u\in\L(E)$. On dit que $u$ est \bf{symétrique} (ou \bf{autoadjoint}) ssi $u=u^*$.\\
    On notera $\m{S}(E)$ l'ensemble des endomorphismes symétriques.
\end{defi}

\begin{defi}{}{}
    Soit $(E,\Lg\.,\.\Rg)$ euclidien et $u\in\L(E)$. On dit que $u$ est \bf{antisymétrique} ssi $u=-u^*$.\\
    On notera $\m{A}(E)$ l'ensemble des endomorphismes antisymétriques.
\end{defi}

\begin{prop}{}{}
    $\m{S}(E)$ et $\m{A}(E)$ sont des $\R$-espace vectoriel.
    \tcblower
    \boxed{\m{S}(E)} C'est un sous-espace vectoriel de $\L(E)$ :\\
    $\hra$ $0\in\m{S}(E)$ et $\forall u,v \in \m{S}(E), ~ \forall \l \in \R, ~ (u+\l v)^* = u^* + \l v^* = u+\l v$.\\
    \boxed{\m{A}(E)} Pareil.
\end{prop}

\begin{prop}{}{}
    Soit $u\in\L(E)$.
    \begin{equation*}
        u\in \m{A}(E) \iff \forall x \in E, ~ \Lg u(x), x \Rg = 0.
    \end{equation*}
    \tcblower
    \boxed{\ra} $\forall x,y \in E, ~ \Lg u(x), y \Rg = -\Lg x,u(y)\Rg$. Puis $y=x$.\\
    \boxed{\la} Soient $x,y\in E$ : $\Lg u(x+y), x+y\Rg = 0$ puis $\Lg u(x), y \Rg = - \Lg x, u(y) \Rg$.
\end{prop}

\begin{prop}{}{}
    Soit $(E,\Lg\.,\.\Rg)$ euclidien et $u\in\L(E)$. Il y a équivalence :
    \begin{enumerate}
        \item $u\in S(E)$.
        \item La matrice de $u$ dans toute b.o.n. de $E$ est symétrique.
        \item Il existe une b.o.n. de $E$ telle que la matrice de $u$ dans cette base soit symétrique.
    \end{enumerate}
    \tcblower
    \boxed{1\ra 2} Soit $\B$ b.o.n. de $E$, alors $\Mat_\B(u)^\top = \Mat_\B(u^*)$ donc $u=u^* \ra \Mat_\B(u)=\Mat_\B(u)^\top$.\\
    \boxed{2\ra3} Trivial.\\
    \boxed{3\ra1} Déjà fait.
\end{prop}

\begin{lemme}{}{}
    \begin{equation*}
        \forall u \in \L(E), ~ \Tr(u^*) = \Tr(u).
    \end{equation*}
    \tcblower
    Soit $\B=(e_1,...,e_n)$ une b.o.n. de $E$. Alors $\Mat_\B(u)=(\Lg u(e_j), e_i \Rg)_{i,j}$.
    \begin{equation*}
        \Tr(u) = \sum_{i=1}^n \Lg u(e_i), e_i \Rg = \sum_{i=1}^n \Lg e_i, u^*(e_i)\Rg = \Tr(u^*).
    \end{equation*}
    Donc $\Tr(u^*v)=\Tr((u^*v)^*)=\Tr(v^*u)$.
\end{lemme}

\begin{exercice}{}{}    \begin{equation*}
        \forall u \in \L(E), ~ \Tr(u^*) = \Tr(u).
    \end{equation*}
    Soit $(E,\Lg\.,\.\Rg)$ euclidien, on pose $\Lg u,v \Rg=\Tr(u^*v)$ pour $u,v\in\L(E)$.\\
    Montrer que l'on a défini un produit scalaire sur $\L(E)$. Montrer que $\L(E)=\S(E)\oplus \A(E)$.
    \tcblower
    Posons $\phi:(u,v)\mapsto \Tr(u^*v)$.\\
    $\hra$ Symétrique : lemme.\\
    $\hra$ Bilinéaire : soient $u,u',v\in\L(E)$ et $\l\in\R$. Alors :
    \begin{equation*}
        \phi(u+\l u', v) = \Tr((u+\l u')^*v) = \Tr((u^* + \l u'^*)v) = \Tr(u^*v + \l u'^*v)=\phi(u,v) + \l \phi(u',v). 
    \end{equation*}
    $\hra$ Positive : $\phi(u,u) = \Tr(u^*u)=\sum_{i=1}^n\Lg (u^*u)(e_i),e_i \Rg = \sum_{i=1}^n\Lg u(e_i), u(e_i)\Rg = \sum_{i=1}^n\|u(e_i)\|^2\geq0$.\\
    $\hra$ Définie : $\phi(u,u) = 0 \ra \sum_{i=1}^n\|u(e_i)\|^2=0 \ra \forall i \in \lb1,n\rb, ~ u(e_i)=0_E$, $u$ est nul sur une base donc nul.\\
    On a déjà $\forall u \in \L(E), ~ u=\frac{u+u^*}{2} + \frac{u-u^*}{2}$ donc $\L(E)=\S(E)+\A(E)$.\\
    Ensuite, pour $u\in\S(E)$ et $v\in\A(E)$, $\Lg u,v\Rg = \Tr(u^*v) = \Tr(uv)$ et $\Lg v,u\Rg = \Tr(v^*u)=-\Tr(vu)=-\Tr(uv)$. 
\end{exercice}

\begin{prop}{}{}
    Soit $(E,\Lg\.,\.\Rg)$ euclidien.
    \begin{enumerate}
        \item Soit $p$ un projecteur de $E$, $p$ est orthogonal ssi $p$ est autoadjoint.
        \item Soit $s$ une symétrie de $E$, $s$ est orthogonale ssi $s$ est autoadjointe.
    \end{enumerate}
    \tcblower
    \boxed{1.\ra} Soit $p$ un projecteur orthogonal sur $F$. Soit $(x,y)\in E^2$.\\
    Alors $\Lg p(x), y \Rg = \Lg p(x), p(y) + (y - p(y)) \Rg = \Lg p(x), p(y)\Rg + 0 = \Lg x, p(y)\Rg$.\\
    \boxed{1.\la} Soit $p$ un projecteur sur $F$ dans la direction $G$. Vérifions $F\bot G$.\\
    Soient $x\in F, y\in G$. Alors $\Lg p(x), y \Rg = \Lg x, p(y)\Rg$ donc $\Lg x,y\Rg = \Lg x, 0_E\Rg = 0$, donc $F\bot G$.\\
    Donc $F\subset G^\bot$, puis $\dim F = \dim G^\bot$ donc $F=G^\bot$.\\
    \boxed{2.} Notons $p$ le projecteur associé à $s$ : $s+\Id_E = 2p$.\\
    Alors $s$ est symétrie orthogonale ssi $G=F^\bot$ ssi $p$ est projection orthogonale ssi $p=p^*$.\\
    Puis $s=2p-Id_E$ donc $s^* = 2p^* - \Id_E = 2p - \Id_E = s$ donc $s$ est autoadjointe.
\end{prop}

\begin{rappel}{}{}
    Soit $(E,\Lg\.,\.\Rg)$ un espace euclidien. Dans une base quelconque, avec $x,y \in E$ et $X,Y$ les vecteurs associés :
    \begin{equation*}
        \Lg x, y \Rg = \sum_{i=1}^n\sum_{j=1}^n x_iy_j\Lg e_i, e_j\Rg = X^\top\Mat(\Lg e_i, e_j \Rg)Y
    \end{equation*}
    Dans une base orthonormée, $\Lg x,y\Rg = X^\top Y = \sum_{i=1}^nx_iy_i$. Et pour $u\in\L(E)$ :
    \begin{equation*}
        \Lg u(x), y \Rg = X^\top\Mat(u)^\top Y.
    \end{equation*}
\end{rappel}

\vspace*{-0.8cm}

\begin{prop}{}{}
    Soit $A\in\M_n(\R)$.
    \begin{enumerate}
        \item $\left[\forall X,Y \in \M_{n,1}(\R), ~ X^\top A Y = 0 \right]\iff A = 0$.
        \item $\left[ \forall X \in \M_{n,1}(\R), ~ X^\top A X = 0 \right] \iff A \in \A_n(\R)$.
    \end{enumerate}
    \tcblower
    \boxed{1.} On spécifie en $E_i$, $E_j$ : $E_i^\top A E_j = a_{i,j} = 0$ pour tout $i,j\in\lb1,n\rb$.\\
    \boxed{2.\ra} $\forall x \in \R^n$, $\Lg x, u(x)\Rg=X^\top A X$ donc $u$ est antisymétrique donc $A$ l'est.\\
    \boxed{2.\la} idem.
\end{prop}

\vspace*{-0.8cm}

\begin{exercice}{}{}
    Soit $(E,\Lg\.,\.\Rg)$ euclidien et $u\in\S(E)$. Montrer que $\Ker u = (\Im u)^\bot$.
    \tcblower
    Soit $x\in\Ker u$, $y\in\Im u$ : $\exists z \in E, ~ y=u(z)$. Alors $\Lg x,y \Rg = \Lg x, u(z) \Rg = \Lg u(x), z \Rg = 0$. Donc $\Ker u \subset (\Im u)^\bot$.\\
    Par théorème du rang, $\dim\Ker u = \dim(\Im u)^\bot$, donc $\Ker u = (\Im u)^\bot$.
\end{exercice}

\vspace*{-0.8cm}

\begin{prop}{Caractérisation des matrices de projection et symétries orthogonales.}{}
    \begin{enumerate}
        \item Soit $P\in\M_n(\R)$. C'est une matrice de projection ssi $P^2=P$.\\
        C'est une matrice de projection orthogonale ssi $P^2=P$ et $P^\top = P$.
        \item Soit $S\in\M_n(\R)$. C'est une matrice de symétrie ssi $S^2=I_n$.\\
        C'est une symétrie orthogonale ssi $S^2=I_n$ et $SS^\top=I_n$ ssi $S^2=I_n$ et $S^\top=S$.
    \end{enumerate}
    \tcblower
    \boxed{1.} On note $p$ l'endomorphisme canoniquement associé à $P$. On a vu $p$ orthogonal ssi $p^2=p$ et $p^*=p$.\\
    Donc $p$ orthogonal ssi $P^2=P$ et $P^\top=P$.\\
    \boxed{2.} On note $s$ l'endomorphisme canoniquement associé à $S$. On a vu $s$ orthogonal ssi $s^2=I_n$ et $s^*=s$.
\end{prop}

\vspace*{-0.8cm}

\begin{prop}{}{}
    Soit $(E,\Lg\.,\.\Rg)$ euclidien, $u\in\S(E)$ et $F$ un sous-espace stable par $u$. Alors $F^\bot$ est stable par $u$.\\
    De plus, les endomorphismes induits $u_{|F}$ et $u_{|F^\bot}$ sont encore symétriques.
    \tcblower
    On rappelle que $F$ est stable par $u$ implique que $F^\bot$ est stable par $u^*$, ici $u=u^*$.
\end{prop}

\begin{ex}{}{}
    Soit $u\in\S(E)$ et $\l,\mu\in\Sp(u)$. Si $\l\neq\mu$, alors $E_\l\bot E_\mu$.
    \tcblower
    Soient $x\in E_\l$ et $y\in E_\mu$. On a $u(x)=\l x$ et $u(y)=\mu y$, et $\Lg u(x), y \Rg = \Lg x, u(y)\Rg$.\\
    Donc $\l \Lg x,y\Rg = \mu \Lg x,y\Rg$ donc $(\l - \mu)\Lg x, y \Rg =0_\R$ donc $\Lg x, y \Rg = 0$ donc $E_\l\bot E_\mu$.
\end{ex}

\begin{thm}{Théorème spectral.}{}
    \begin{enumerate}
        \item Soit $A\in\S_n(\R)$. Alors $A$ est orthogonalement diagonalisable : $\exists P \in O_n(\R), ~ \exists D \in \D_n(\R), ~ A=PDP^{\top}$.
        \item Soit $(E,\Lg\.,\.\Rg)$ euclidien et $u\in\S(E)$. Alors $\Sp(u)\subset\R$, $u$ est dz dans une b.o.n. et $E\bigoplus_{\l\in\Sp(u)}E_\l$.
    \end{enumerate}
    Il est clair que $1\iff 2$.
    \tcblower
    Soit $u\in\S(E)$ et $\l\in\Sp(u)$. Soit $\B$ une b.o.n. de $E$ et $A=\Mat_\B(u)$ : $A\in\S_n(\R)$. Soit $x\in E$ de vecteur $X$.\\
    Alors $u(x) = \l x \iff AX=\l X$. Montrons que $\l\in \R$ : 
    \begin{equation*}
        \ov{X}^\top A X = \ov{X}^\top \l X = \l \ov{X}^\top X = \l \sum_{i=1}^n|x_i|^2.
    \end{equation*}
    Or $\ov{X}^\top AX = (A\ov{X})^\top X = \ov{\l}\ov{X}^\top X = \ov{\l} \sum_{i=1}^n |x_i|^2$ donc $\l=\ov{\l}$ donc $\l\in\R$.\\
    Notons $F=\bigoplus_{\l \in \Sp(u)}E_\l$, c'est un ssev de $E$ stable par $u$, on suppose par l'absurde que $F\subsetneq E$.\\
    Alors $E=F\oplus F^\bot$ et $F^\bot$ est stable par $u$ car symétrique, et $u_{|F^\bot}$ est symétrique : ses valeurs propres sont réelles.\\
    Soit $\mu$ une vp de $u_{|F^\bot}$, et $x$ $\vp$ associé, alors $x\in F^\bot\cap F$ donc $x=0_E$, absurde. Donc $F=E$.\\
    Donc $u$ est diagonalisable par concaténation des b.o.n. de $E_\l$.
\end{thm}

\begin{ex}{}{}
    Orthodiagonaliser la matrice :
    \begin{equation*}
        A=\begin{pmatrix}
          2 & 1 & 0 \\ 1 & 2 & 1 \\ 0 & 1 & 2   
        \end{pmatrix} \in \S_3(\R)
    \end{equation*}
    \tcblower
    D'après le théorème spectral, $A$ est orthogonalement diagonalisable.\\
    On a $\X_A = (\l-2)\left[ (\l - 2)^2 - 2 \right] = (X - 2)(X - 2 - \sqrt{2})(X - 2 + \sqrt{2})$.\\
    $E_2=\Vect(1,0,1)$, $E_{2-\sqrt{2}}=\Vect(1,-\sqrt{2},1)$ et $E_{2+\sqrt{2}}=\Vect(1,\sqrt{2},1)$.
    \begin{equation*}
        P = \begin{pmatrix}\frac{1}{\sqrt{2}} & \frac{1}{2} & \frac{1}{2} \\ 0 & \frac{\sqrt{2}}{2} & -\frac{\sqrt{2}}{2} \\ -\frac{1}{\sqrt{2}} & \frac{1}{2} & \frac{1}{2}\end{pmatrix}\in O_3(\R)
    \end{equation*}
    Puis $A=PDP^\top$ avec $D=\Diag(2,2+\sqrt{2},2-\sqrt{2})$.
\end{ex}

\begin{ex}{}{}
    Orthodiagonaliser la matrice :
    \begin{equation*}
        A = \begin{pmatrix}1&1&1\\1&1&1\\1&1&1\end{pmatrix} \in \S_3(\R)
    \end{equation*}
    \tcblower
    Par théorème spectral, elle est orthogonalement diagonalisable.\\
    On a $\rg(A)=1$, donc $\dim\Ker A = 2$, donc $0\in\Sp(A)$ et $\dim E_0 = 2$ et $\Tr(A)=3$ donc $3\in\Sp(A)$.\\
    On a $E_3=\Vect((\frac{1}{\sqrt{3}}, \frac{1}{\sqrt{3}}, \frac{1}{\sqrt{3}}))$ et $E_0=E_3^\bot=\Vect((-1,1,0),(1,1,-2))$...
\end{ex}

\begin{ex}{}{}
    \begin{equation*}
        \begin{pmatrix}
            i&1\\1&-i 
        \end{pmatrix} \in \S_3(C)
    \end{equation*}
    \tcblower
    Le thérorème spetral ne s'applique pas ici car $A\notin\S_2(\R)$.\\
    On a $\X_A=X^2-\Tr(A)X+\det(A)=X^2$, donc $\Sp(A)=\{0\}$, donc $A$ est nilpotente, pas diagonalisable car $A\neq0$.
\end{ex}

\begin{prop}{}{}
    \begin{enumerate}
        \item Soit $A,B\in\S_n(\R)$ tels que $AB=BA$. Alors $A$ et $B$ sont orthogonalement diagonalisables.
        \item Pour $u,v\in\S(E)$, si $uv=vu$, alors il existe une b.o.n. commune de $\vp$ de $u$ et $v$.
    \end{enumerate}
    \tcblower
    Soit $u\in \S(E)$. Alors $E=\bigoplus E_{\l}$, et $uv=vu$ implique $\forall \l \in \Sp(u)$, $v(E_\l)\subset E_\l$.\\
    Notons $v_\l = v_{|E_\l}\in \L(E_\l)$, c'est encore un endomorphisme symétrique : $v_\l\in\S(E_\l)$.\\
    D'après le théorème spectral, on peut choisir une b.o.n. $\B_\l$ de $E_\l$ constituée de $\vp$ de $v_\l$.\\
    On pose $\B$ la concaténation des $\B_\l$, c'est une b.o.n. de $E$ et les vecteurs de $\B$ sont des $\vp$ de $v$, mais aussi de $u$.
\end{prop}

\begin{prop}{Réciproque du théorème spectral.}{}
    \begin{enumerate}
        \item Soit $A\in\M_n(\R)$ telle que $A=PDP^{-1}$ avec $P\in O_n(\R)$ et $D\in O_n(\R)$, alors $A\in\S_n(\R)$.
        \item Soit $u\in\L(E)$, si il existe $\B$ b.o.n. de $E$ constituée de $\vp$ de $u$, alors $u\in\S(E)$.
    \end{enumerate}
    \tcblower
    \boxed{1.} $A=PDP^{\top}$ donc $A^\top = PD^\top P^\top = A$.\\
    \boxed{2.} Soit $\B=(e_1,...,e_n)$ une b.o.n. de $E$ telle que $\forall i \in \lb1,n\rb, ~ u(e_i)=\l e_i$.\\
    Soient $x,y \in E$. $\Lg u(x), y \Rg = \sum_{i=1}^n y_iu(x_i)=\sum_{i=1}^n y_i(\l_i x_i)$ et $\Lg x,u(y)\Rg = \sum_{i=1}^n x_i(\l_i y_i)$.
\end{prop}

\pagebreak

\begin{exercice}{Rayon spectral.}{}
    \begin{enumerate}
        \item Soit $E$ un evdf, $u\in\L(E)$, $p(u)=\max\{|\l|, ~ \l \in \Sp(u)\}$ (rayon spectral). Montrer que $p(u)\leq\abs u \abs$.
        \item Pour $(E,\Lg\.,\.\Rg)$ un espace euclidien, $u\in\S(E)$, $\a=\min\{\l,\l\in\Sp(u)\}$ et $p=\max\{\l,\l\in\Sp(u)\}$.\\
        Montrer que $\forall x \in E, ~ \a\|x\|^2 \leq \Lg u(x), x \Rg \leq \b \|x\|^2$.
        \item Soit $(E,\Lg\.,\.\Rg)$ euclidien et $u\in\S(E)$. Montrer que $p(x)=\abs u\abs$.
    \end{enumerate}
    \tcblower
    \boxed{1.} Supposons que $u$ admet au moins une valeur propre réelle.\\
    Soit $\l$ une valeur propre de $u$ : $\exists x \in E, ~ u(x)=\l x$ puis $\|u(x)\|=|\l|\|x\|$ puis $|\l|=\frac{\|u(x)\|}{\|x\|}\leq\abs u\abs$.\\
    Par passage au max : $p(u)\leq\abs u\abs$.\\
    \boxed{2.} Soit $u\in S(E)$, d'après le théorème spectral, il existe $\B=(e_1,...,e_n)$ b.o.n. de $E$ de $\vp$ de $u$ : $u(e_i)=\l_ie_i$.\\
    Alors $\forall i \in \lb1,n\rb, ~ \exists \l_i \in \R, ~ u(e_i)=\l_i e_i$. Soit $x\in E$ : $u(x)=\sum_{i=1}^n(\l_i x_i)e_i$.\\
    Puis $\Lg u(x), x \Rg = \sum_{i=1}^n\l_i x_i^2$ et $\|x\|^2 = \sum_{i=1}^nx_i^2$ car $\B$ est une b.o.n.\\
    Par choix de $\a$ et $\b$ : $\forall i \in \lb1,n\rb, ~ \a\leq \l_i \leq \b$ puis $\a x_i^2\leq \l_ix_i^2 \leq \b x_i^2$.\\
    Donc $\a\|x\|^2\leq\Lg u(x), x \Rg \leq \b\|x\|^2$ en sommant.\\
    \boxed{3.} On a déjà $p(u)\leq\abs u\abs$. Soit $x\in E$ et $\B=(e_1,...,e_n)$ une b.o.n. de $\vp$ de $u$.\\
    Alors $\|u(x)\|^2=\sum_{i=1}^n(\l_ix_i)^2$ et $\|x\|^2=\sum_{i=1}^nx_i^2$ donc $\|u(x)\|^2\leq p(u)^2\|x\|^2$ donc $\forall x \neq 0, ~ \frac{\|u(x)\|}{\|x\|}\leq p(u)$.\\
    Par passage au sup, $\abs u \abs \leq p(u)$. Donc $p(u)=\abs u\abs$.
\end{exercice}

\section{Endomorphismes symétriques positifs.}

\begin{defi}{Endomorphisme symétrique positif.}{}
    Soit $A\in\S_n(\R)$.
    \begin{enumerate}
        \item $A$ est positive ssi $\forall X \in \M_{n,1}(\R), ~ X^\top AX \geq 0$.
        \item $A$ est définie positive ssi $\forall X \in \M_{n,1}(\R)$ non nul, $X^\top AX>0$.
    \end{enumerate}
\end{defi}

\begin{defi}{}{}
    Soit $(E,\Lg\.,\.\Rg)$ un espace euclidien, $u\in\S(E)$.
    \begin{enumerate}
        \item $u$ est positive ssi $\forall x \in E, ~ \Lg u(x), x \Rg \geq 0$.
        \item $u$ est strictement positive ssi $\forall x \in E, ~ x\neq0 \ra \Lg u(x), x \Rg > 0$.
    \end{enumerate}
\end{defi}

\begin{prop}{}{}
    Soit $(E,\Lg\.,\.\Rg)$ euclidien et $u\in\S(E)$. Il y a équivalence :
    \begin{enumerate}
        \item $u$ est positif.
        \item La matrice de $u$ dans n'importe quelle base de $E$ est positive.
        \item Il existe une b.o.n. de $E$ telle que la matrice de $u$ dans $E$ soit positive.
    \end{enumerate}
    \tcblower
    \boxed{1\ra2} Soit $\B=(e_1,...,e_n)$ une b.o.n. de $E$ et $A=\Mat_\B(u)$.\\
    Déjà, $u\in\S(E)$ donc $A\in\S_n(\R)$ et pour $X=(x_i)_i\in\M_{n,1}(\R)$ avec $x=\sum_{i=1}^n x_ie_i$.\\
    On sait que $\Lg x, u(x)\Rg = X^\top A X$ donc $X^\top AX\geq 0$.\\
    \boxed{2\ra3} trivial; \boxed{3\ra 1} facile.
\end{prop}

\begin{nota}{}{}
    On notera $\S_n^+(\R)$ l'ensemble des matrices symétriques réelles positives.\\
    On notera $\S_n^{++}(\R)$ l'ensemble des matrices symétriques réelles définies positives.\\
    On notera $\S^+(E)$ l'ensemble des matrices symétriques de $E$ positives.\\
    On notera $\S^{++}(E)$ l'ensemble des matrices symétriques de $E$ définies positives. 
\end{nota}

\begin{prop}{Propriétés topologiques.}{}
    \begin{enumerate}
        \item $\S_n(\R)$ est un fermé de $\M_n(\R)$.
        \item $\S_n^+(\R)$ est un fermé de $\M_n(\R)$.
    \end{enumerate}
    \tcblower
    \boxed{1.} Soit $(A_k)_k\in\S_n(\R)^\N$ convergente de limite $A$. Par hypothèse, $\forall k \in \N, ~ A_k^\top = A_k$.\\
    De plus, la transposition est linéaire en dimension finie, donc continue sur $\M_n(\R)$.\\
    Par passage à la limite, $A^\top = A$, donc $A\in\S_n(\R)$.\\
    \boxed{2.} Soit $(A_k)_k\in\S_n^+(\R)^\N$ convergente de limite $A$. Déjà, $A\in\S_n(\R)$ par la question précédente.\\
    Ensuite, pour $X\in\M_{n,1}(\R)$, on a $\forall k \in \N, ~ X^\top A_k X \geq 0$.\\
    Puis $M\mapsto X^\top M X$ est continue sur $\M_n(\R)$ car lin. en dimension finie, donc par passage à la limite, $X^\top AX \geq 0$.\\
    C'est vrai pour tout $X\in\M_{n,1}(\R)$, donc $A\in\S_n^+(\R)$. 
\end{prop}

\begin{exercice}{}{}
    Montrer que si $A\in\S_n^+(\R)$, alors $\forall i \in \lb1,n\rb, ~ [A]_{i,i}\geq0$.
    \tcblower
    Soit $A=(a_{i,j})_{i,j}\in\S_n^+(\R)$.\\
    Par hypothèse, $\forall X \in \M_{n,1}(\R)$, $X^\top AX\geq0$. On spécifie en $E_{i_0}$ :
    \begin{equation*}
        \forall i_0 \in \lb1,n\rb, ~ E_{i_0}^\top A E_{i_0} = a_{i_0,i_0}\geq 0.
    \end{equation*}
    \bf{Remarque.} Si $A\in\S_n^{++}(\R)$, alors $\forall i \in \lb1,n\rb, ~ [A]_{i,i}>0$.
\end{exercice}

\begin{prop}{Caractérisation des matrices symétriques positives.}{}
    \begin{enumerate}
        \item $A \in \S_n^{+}(\R) \iff \Sp A \subset \R_+$
        \item $A \in \S_n^{++}(\R) \iff \Sp A \subset \R_+^*$ (facile avec 1.)
        \item $\S_n^{++}(\R)=\S_n^+(\R)\cap\GL_n(\R)$.
    \end{enumerate}
    \tcblower
    \boxed{1.\ra} Soit $\l\in\Sp(A)$ : $\exists X\in\M_{n,1}(\R), ~ AX=\l X$. Puis $X^\top AX = \l X^\top X = \l \geq 0$ car $A\in\S_n^+(\R)$.\\
    \boxed{1.\la} $A$ est alors symétrique réelle puis par théorème spectral : $A$ est orthogonalement diagonalisable.\\
    Ainsi, $\exists P \in O_n(\R), ~ \exists D \in \D_n(\R), ~ A=PDP^\top$.\\
    Alors pour $X\in \M_{n,1}(\R), ~ X^\top AX = X^\top PDP^\top X$, on pose $Y=P^{\top}X\in\M_{n,1}(\R)$.\\
    Alors $X^\top AX = Y^\top DY = \sum_{i=1}^n\l_i y_i^2 \geq 0$, donc $A\in\S_n^+(\R)$.\\
    \boxed{3.} On rappelle que $A\in\GL_n(\R) \iff 0 \notin \Sp(A)$, puis $A\in\S_n^{++}(\R) \iff \Sp(A) \subset \R_+^*$ ssi $A\in\S_n^+(\R)\cap \GL_n(\R)$
\end{prop}

\begin{exercice}{}{}
    \begin{enumerate}
        \item Soit $A\in\M_n(\R)$. Vérifier que $AA^\top\in\S_n^+(\R)$ et $A^\top A \in \S_n^+(\R)$.
        \item Soit $A\in\GL_n(\R)$. Vérifier que $AA^\top\in\S_n^{++}(\R)$ et $A^\top A\in\S_n^{++}(\R)$.
    \end{enumerate}
    \tcblower
    \boxed{1.} Déjà, $(AA^\top)^\top = (A^\top)^\top A^\top = AA^\top$ donc $AA^\top \in \S_n(\R)$; de même, $A^\top A \in \S_n(\R)$.\\
    Pour $X\in\M_{n,1}(\R), ~ X^\top(AA^\top)X = (X^\top A)A^\top X = (A^\top X)^\top A^\top X = \|A^\top X\|^2 \geq 0$.\\
    \boxed{2.} Pour $A\in\GL_n(\R)$, $\det(AA^\top)=\det A \det A^\top = \det(A)^2 \neq 0$ : $AA^\top \in \GL_n(\R)$, idem pour $A^\top A$.\\
    Puis $AA^\top \in \S_n^+(\R)$ avec 1, donc $AA^\top \in \S_n^+(\R)\cap \GL_n(\R)=\S_n^{++}(\R)$.
\end{exercice}

\begin{exercice}{Matrice de Hilbert.}{}
    Montrer que les valeurs propres de $H=\left( \frac{1}{i+j-1} \right)_{i,j}\in\M_n(\R)$ sont toutes positives.
    \tcblower
    \bf{Ne pas chercher à calculer $\X_H$.}
    \begin{equation*}
        H = \begin{pmatrix} 1 & \frac{1}{2} & ... & \frac{1}{n} \\ \frac{1}{2} & \frac{1}{3} & ... & \frac{1}{n+1} \\ \vdots & \vdots & & \vdots \\ \frac{1}{n} & \frac{1}{n+1} & ... & \frac{1}{2n-1}\end{pmatrix}
    \end{equation*}
    Déjà, $H$ est symétrique réelle, donc par théorème spectral, les vp de $H$ sont toutes réelles et $H$ est orthodz.\\
    On sait que $H\in\S_n^+(\R) \iff \Sp(H)\subset\R_+$. Soit $X=(x_i)\in\M_{n,1}(\R)$. On remarque $\int_0^1 t^k\dt = \frac{1}{k+1}$.
    \begin{align*}
        X^\top HX = \sum_{i,j} x_ix_j\frac{1}{i+j-1} = \sum_{i,j}x_ix_j \int_0^1t^{i+j-2}\dt = \int_0^1 \sum_{i,j} x_i x_j t^{i-1}t^{j-1}\dt = \int_0^1 \left( \sum_{i=1}^n x_i t^{i-1} \right)^2\dt \geq 0.
    \end{align*}
    Donc $\forall X \in \M_{n,1}(\R), X^\top H X \geq 0$ donc $H \in \S_n^+(\R)$. Comme $H\in\S_n^+(\R), ~ H \in \GL_n(\R) \iff H\in\S_n^{++}(\R)$.\\
    Supposons $X^\top H X = 0$. Montrons que $X=0$. On a $\int_0^1\left( \sum_{i=1}^n x_it^{i-1} \right)\dt=0$.\\
    C'est une intégrale nulle d'une fonction \bf{continue} positive, donc $\forall t \in [0,1], ~ \sum_{i=1}^n x_it^{i-1} = 0$.\\
    C'est un polynôme nul sur $[0,1]$, il est nul partout par rigidité. Donc $\forall i \in \lb1,n\rb, ~ x_i = 0$ donc $X=0$.\\
    Ainsi, $H\in\GL_n(\R)$ donc $H\in\S_n^{++}(\R)$.
\end{exercice}

\newcommand*{\Gram}{\nt{Gram}}

\begin{exercice}{}{}
    Soit $(E,\Lg\.,\.\Rg)$ préhilbertien réel et $(a_1,...,a_n) \in E^n$.
    \begin{enumerate}
        \item Montrer que $\Gram(a_1,...,a_n)\in\S_n^+(\R)$.
        \item Montrer que $\Gram(a_1,...,a_n)\in\S_n^{++}(\R) \iff (a_1,...,a_n)$ libre.
        \item Retrouver que $H\in\S_n^{++}(\R)$.
    \end{enumerate}
    \tcblower
    \boxed{1.} Déjà, $\Gram(a_1,...,a_n)\in\S_n(\R)$ par symétrie du produit scalaire. Soit $X=(x_i)\in\M_{n,1}(\R)$.\\
    On a $X^\top\Gram(a_1,...,a_n)X = \sum_{i=1}^n x_i x_j \Lg a_i, a_j \Rg = \Lg \sum_{i=1}^n x_i a_i, \sum_{j=1}^n x_j a_j \Rg = \|\sum_{i=1}^n x_i a_i\|^2 \geq 0$.\\
    \boxed{2.} $\Gram(a_1,...,a_n)\in\GL_n(\R)$ ssi $(a_1,...,a_n)$ est libre et $\S_n^{++}(\R)=\GL_n(\R)\cap \S_n^{+}(\R)$ donc...\\
    \boxed{3.} On pose $E=\R[X]$ muni de $\Lg P,A\Rg = \int_0^1 P(t)Q(t)\dt$. On pose $(P_1,...,P_n)$ où $P_i=X^{i-1} ~ : ~ \Lg P_i, P_j \Rg = \frac{1}{i+j-1}$.\\
    Alors $H=\Gram(P_1,...,P_n)$ puis $H\in\S_n^{++}(\R)$ car $(1,X,...,X^{n-1})$ est libre.
\end{exercice}

\begin{exercice}{}{}
    Soit $A\in\S_n^+(\R)$. Est-elle une matrice de Gram ?
    \tcblower
    On munit $\R^n$ du produit scalaire usuel.\\
    Par théorème spectral : $\exists P \in O_n(\R), ~ \exists D \in \D_n(\R), ~ A=PDP^\top$ où $D=\Diag(\l_i)$.\\
    Soit $\Delta=\Diag(\sqrt{\l_i})$, alors $A=P\Delta^2P^\top = P\D(P\Delta)^\top$ et pour $M=(P\Delta)^\top, ~ A=M^\top M$.\\
    Notons $v_1,...,v_n$ les colonnes de $M$. Alors $[M^\top M]_{i,j}=\Lg v_i, v_j\Rg$ donc $A=(\Lg v_i, v_j\Rg)$.\\
    Donc $A=\Gram(v_1,...,v_n)$.
\end{exercice}

\vspace*{-0.8cm}

\begin{exercice}{Inégalité d'Hadamard.}{}
    Soit $\R^n$ euclidien canonique et $P\in\GL_n(\R)$.
    \begin{enumerate}
        \item Soit $\B$ b.o.n. de $\R^n$ telle que $P=\Mat_{\B_c,\B}$.\\
        On se donne $\cursive{C}$ une orthonormalisée de Schmidt de $\B$ et $Q=\Mat_{\B,\cursive{C}}$.\\
        Montrer que $PQ \in O_n(\R)$, en déduire $P=OT$ avec $O\in O_n(\R)$ et $T\in\T_n^{++}(\R)$ avec unicité.
        \item Soit $A\in\M_n(\R)$. Montre que $\det(A)^2 \leq \prod_{i=1}^n \left( \sum_{j=1}^n a_{i,j}^2 \right)$.
    \end{enumerate}
    \tcblower
    \boxed{1.} $PQ$ est la matrice de passage de $\B_C$ à $\cursive{C}$, donc $PQ\in\O_n(\R)$ et $P=PQQ^{-1}$. On pose $T=Q^{-1}$.\\
    Notons $\cursive{C}=(a_1,...,a_n)$ ses colonnes et $\B=(e'_1,...,e'_n)$, alors $Q=(\Lg e'_j, e_i\Rg)_{i,j}$.\\
    De plus, $\cursive{C}$ est orthonormalisée de Schmidt : $\forall i \in \lb1,n\rb, ~ e'_i \in \Vect(e_1,...,e_n)$ et $\Lg e'_i,e_i\Rg>0$.\\
    Donc $Q^{-1}$ est triangulaire supérieure réelle à diagonale strictement positive.\\
    Sq $P=OT=O'T'$, avec $O,O'\in O_n(\R)$, $T,T'\in\T_n^{++}(\R)$ : $(O')^{-1}O=T'T^{-1}$ or $O_n(\R)\cap\T_n^{++}(\R)=\{ I_n\}$ donc $O'=O$ et $T=T'$.\\
    \boxed{2.} Si $A\notin \GL_n(\R)$, c'est trivial.\\
    Si $A=OT$, où $O\in O_n(\R)$ et $T\in \T_n^{++}(\R)$, on note $e'_1,...,e'_n$ les colonnes de $A$. C'est une base car $A\in\GL_n(\R)$.\\
    On note $e_1,...,e_n$ son orthonormalisée de Schmidt, alors \begin{equation*}\det(A)^2=\det(O)^2\det(T)^2=\det(T)^2=\prod_{i=1}^n\Lg e'_i, e_i\Rg^2 \leq \prod_{i=1}^n \|e'_i\|^2\cancel{\|e_i\|^2}=\prod_{i=1}\sum_{j=1}^n a_{j,i}^2\end{equation*}
\end{exercice}

\vspace*{-0.8cm}

\begin{exercice}{Décomposition polaire.}{}
    \begin{enumerate}
        \item Soit $A\in\GL_n(\R)$. Montrer que $A^\top A\in\S_n^{++}(\R)$.\\
        En déduire que $\exists!(O,S)\in O_n(\R)\times S\in\S_n^{++}(\R)$ tel que $A=OS$.
        \item Soit $A\in\M_n(\R)$. Montrer que $\exists (O,S)\in O_n(\R)\times \S_n^+(\R), ~ A=OS$. 
    \end{enumerate}
    \tcblower 
    \boxed{1.} On sait déjà que $A^\top A \in \S_n^{++}(\R)$. Supposons qu'il existe $O,S$ tels que $A=OS$.\\
    Alors $A^\top A = S^\top O^\top O S = S^2$, or toute matrice de $\S_n^{++}(\R)$ possède une unique racine carrée dans $\S_n^{++}(\R)$ donc $S$ est déterminée de manière unique puis $O=AS^{-1}$ aussi.\\
    Soit alors $S$ l'unique racine de $A^\top A$ dans $\S_n^{++}(\R)$ et $O=AS^{-1}$, alors $O=AS^{-1}\ra OS=A$, $OO^\top = I_n$.\\
    \boxed{2.} Soit $A\in\M_n(\R)$. Par densité de $\GL_n(\R)$ dans $\M_n(\R)$, $\exists (A_n)\in\GL_n(\R)^\N$ telle que $A_n\to A$.\\
    On a $\forall k \in \N, ~ \exists O_k \in O_n(\R), ~ \exists S_k \in \S_n^{+}(\R), ~ A_k = O_kS_k$.\\
    Or $O_n(\R)$ est compact, donc $\exists \phi : \N\to \N$ strictement croissante telle que $O_{\phi(k)}\to O\in O_n(\R)$.\\
    On a $A_{\phi(k)}\to A$ (extraite), donc $S_{\phi(k)}=O^\top_{\phi(k)}A_{\phi(k)}\to O^\top A$ par continuité de la transposée et du produit.\\
    Alors $A=OS$ avec $O\in O_n(\R)$ et $S_n^{+}(\R)$ est fermé donc $S=O^\top A\in\S_n^{+}(\R)$.
\end{exercice}

\begin{exercice}{Racine carrée de matrices symétriques positives.}{}
    \begin{enumerate}
        \item Soit $A\in\S_n^+(\R)$. Montrer que $\exists B \in \S_n^+(\R), ~ B^2=A$.
        \item Soit $A\in\S_n^{++}(\R)$. Montrer que $\exists B \in \S_n^{++}(\R)$ tel que $B^2=A$.
        \item Montrer l'unicité de $B$.
    \end{enumerate}
    \tcblower
    \boxed{1.} $A$ est symétrique réelle donc par théorème spectral : $\exists P \in O_n(\R), ~ \exists D \in \D_n(\R), ~ A=PDP^\top$.\\
    De plus, $A\in\S_n^+(\R)$ donc $D=\Diag(\l_i)$ avec $\forall i \in \lb1,n\rb,~\l_i \in \Sp(A) \et \l_i\geq0$.\\
    On pose $\Delta = \Diag(\sqrt{\l_i})$ et $B=P\Delta P^{-1}$, alors $B^2=P\Delta^2 P^{-1}=A$, $B^\top = B$ et $\Sp(B)\subset\R_+$.\\
    \boxed{2.} $A\in\S_{n}^{++}(\R)$ donc $A\in\GL_n(\R)$ donc $\l_i\neq0$ donc $\Sp(B)\subset\R_+^*$.\\
    \boxed{3.} Soit $A\in\S_n^+(\R)$ et $B\in\S_n^+(\R)$ tels que $A=B^2$.\\
    On a $AB=B^2B=BB^2=BA$ puis $A,B\in\S_n^+(\R)$ : $A=PDP^\top$ et $B=P\Delta P^\top$.\\
    Puis $B^2=A \iff \Delta^2= D \iff \forall i \in \lb1,n\rb, ~ \mu_i^2=\l_i \iff \forall i \in \lb1,n\rb, ~ \mu_i = \sqrt{\l_i}$.\\
    Avec $\Sp(A)=\{\l_i\}$ et $\Sp(B)=\{\mu_i\}$. On a donc montré que $B$ est determiné de manière unique.
\end{exercice}

\begin{exercice}{}{}
    Soit $(E,\Lg\.,\.\Rg)$ euclidien et $u\in\S^+(E)$.
    \begin{enumerate}
        \item Montrer qu'il existe $v\in\S^+(E)$ tel que $v^2=u$.
        \item Montrer l'unicité de $v$.
    \end{enumerate}
    \tcblower
    \boxed{1.} Par le théorème spectral, il existe $\B=(e_1,...,e_n)$ base de $E$ de $\vp$ de $u$.\\
    On pose $v$ tel que $\forall i \in \lb1,n\rb, ~ v(e_i) = \sqrt{\l_i}e_i$. Alors $v^2(e_i)=\l_ie_i=u(e_i)$ donc $v^2=u$.\\
    De plus, $v\in\S^+(\R)$ par réciproque du théorème spectral.\\
    \boxed{2.} Soit $v\in\S^+(E)$ tel que $v^2=u$. On reprend $\B=(e_1,...,e_n)$ b.o.n. de $E$.\\
    Notons $E_{\l_i}=\Ker(u-\l_i\Id_E)$, alors $E=\bigoplus_{\l\in\Sp(u)}E_\l$ et $\forall \l \in \Sp(u), ~ v(E_\l)\subset E_\l$.\\
    Notons $v_i=v_{|E_{\l_i}}$ induit et diagonalisable. On a $\forall x \in E_{\l_i}, ~ v_i(x)=\l_i x$ alors $X^2-\l_i$ annule $v_i$.\\
    On a donc $\Sp(v_i)\subset\{-\sqrt{\l_i},\sqrt{\l_i}\}$. Or $v_i\in\S^+(E)$ donc $\Sp(v_i)=\{\sqrt{\l_i}\}$. Or $v_i$ dz donc $v_i=\sqrt{\l_i}\Id_E$.\\
    Donc tout $v_i$ est determiné uniquement donc $v$ aussi.
\end{exercice}

\begin{exercice}{}{}
    Soit $A\in\A_n(\R)$ et $f$ canoniquement associé avec $n\in\N$ impair.
    \begin{enumerate}
        \item Montrer que $\forall (x,y) \in E^2, ~ \Lg f(x),y\Rg = -\Lg x,f(y)\Rg$.
        \item Montrer que $f$ admet 0 comme unique valeur propre réelle.
        \item Montrer que $I_n-A\in\GL_n(\R)$ et $I_n+A\in\GL_n(\R)$.
        \item Soit $\B=(I_n-A)^{-1}(I_n+A)$. Montrer que $B\in\SO_n(\R)$.
    \end{enumerate}
    \tcblower
    \boxed{1.} Pour $X,Y\in\M_{n,1}(\R), ~ \Lg AX,Y \Rg = (AX)^\top Y = X^\top (-A)Y = -\Lg X,AY \Rg$.\\
    \boxed{2.} Soit $\l\in\Sp(f)$ : $\exists x \in E^*, ~ f(x) = \l x$. Alors $\Lg f(x),x\Rg = -\Lg x,f(x)\Rg$.\\
    Alors $\Lg \l x, x \Rg = -\Lg x, \l x \Rg$ puis $\l\|x\|^2 = -\l\|x\|^2$ donc $\l=0$ donc $\Sp(f)\subset\{0\}$.\\
    On a $\det(A)=\det(A^\top)=(-1)^n\det(A)=-\det(A)=0$ donc $0$ est vp. Donc $\Sp(f)=\{0\}$.\\
    \boxed{3.} $\X_A(\l)=\det(\l I_n-A)$ et $\l I_n-A\in\GL_n(\R) \iff \X_A(\l)\neq0 \iff \l \notin \Sp(A)$.\\
    Or $1\notin\Sp(A)$ donc $I_n-A\in\GL_n(\R)$ puis $(I_n-A)^\top=I_n+A$ donc $I_n+A\in\GL_n(\R)$.\\
    \boxed{4.} $B^\top B=...=I_n$ et $\det(B)=\frac{\det(I_n+A)}{\det(I_n-A)}=\frac{\det(*^\top)}{\det(*)}=1$.
\end{exercice}

\end{document}
